\section{Experiments}
\label{sec:experiments}

% Mirrors Wu & Tu Section 4: three subsections (Tasks and Datasets / Settings / Results), about
% 1.6 pages, and -- the part worth copying -- ONE results table for the whole paper. Their
% Table 1 covers five tasks and six datasets in eleven rows, and their Results prose is barely
% half a column. Detailed dataset descriptions go to the appendix.

\NOTE{No number appears here without a row in the run log of the corresponding
\texttt{experiments/*/EXPERIMENT\_STATUS.md}: date, commit, config, seed, metric, wall-clock.}

\subsection{Tasks and Datasets}
\label{sec:tasks}

\TODO{brief, with the detail in \Cref{app:datasets} --- this is how Wu \& Tu handle it.}

\paragraph{Language modelling.} \TODO{PTB \citep{marcus1993building} and/or WikiText-2
\citep{merity2017pointer}; perplexity on the held-out set.}

We deliberately do not use WikiText-103. \citet{wu2023probabilistic} report that the encoder
model significantly underperforms transformers beyond roughly $10^5$ sentences, suspected to be
the absence of a feed-forward structure; running there would test a known failure mode of the
encoder rather than the causal extension. \NOTE{state this before a reviewer asks why the
corpus is small --- their own paper makes the same point in its Discussion}

\subsection{Settings}
\label{sec:settings}

\TODO{parameter budget 20--50M matched across models; identical tokenizer, vocabulary and
context length; one shared training loop; Adam with identical hyperparameters and schedule;
$k$ seeds, mean $\pm$ std. Wu \& Tu report mean and std over 5 random runs --- match that.}

\paragraph{Matching budgets.}
We report parameter count \emph{and} wall-clock, since the causal PT shares parameters across
iterations and equal parameter count therefore does not mean equal compute.

\paragraph{The embedding split.}
We report embedding and non-embedding parameters separately, which matters more here than in a
usual comparison. With tying forced (\Cref{sec:output}) the causal PT's parameters sit almost
entirely in $\Sfac \in \R^{|\Vocab|\times\nlab}$; the arc scores $\Tfac{c}$, the word unary
$\bfac$ and the root key $\rfac{c}$ are negligible beside it. At a matched total the causal PT
is close to an embedding table plus a little and a transformer baseline is not, and a single
total hides that completely. \TODO{state the matching rule chosen --- total parameters,
non-embedding parameters, or both with vocabulary fixed. Any is defensible; leaving it unstated
is what gets attacked first.}

\paragraph{Loop depth.}
The loop depth of the Looped baseline matches the number of MFVI iterations $\niter$, so
effective depth is comparable as well as parameter count. We report a sweep over $\niter$ rather
than a single value, so the comparison does not rest on one lucky setting.

\subsection{Results}
\label{sec:results}

\begin{table*}[t]
\centering
\small
\begin{tabular}{llcccc}
\toprule
Dataset & Metric & \makecell{Params\\(emb.\ / non-emb.)} & Transformer & Looped & Causal PT \\
\midrule
PTB         & Perplexity & \TODO{} & \TODO{} & \TODO{} & \TODO{} \\
WikiText-2  & Perplexity & \TODO{} & \TODO{} & \TODO{} & \TODO{} \\
\bottomrule
\end{tabular}
\caption{Main results, mean $\pm$ std over \TODO{$k$} random runs. The comparison this paper is
about is the last two columns.}
\label{tab:main}
\end{table*}

\paragraph{Correctness of the implementation.}
\label{sec:exp-correctness}
\TODO{two sentences plus \Cref{app:validation}. Only the two checks that test the derivation
rather than the code belong in the main text: the mean-field free energy is non-increasing
across iterations on a toy graph, and \Cref{eq:exact-readout} agrees with brute-force
enumeration over a tiny vocabulary to \TODO{1e-N}. Causality, tying, normalisation and the
single-batch overfit are software checks --- list them in the appendix.}

\paragraph{Language modelling.}
\TODO{read \Cref{tab:main}. Two readings, in this order: (i) the causal PT trains and lands
within a modest gap of the transformer --- necessary but not the point; (ii) the causal PT
vs.\ Looped gap is the paper's claim, because Looped has the weight sharing without the
structure (\Cref{tab:three-point}). If the causal PT is worse than Looped across the $\niter$
sweep, the central claim fails and we say so here.}

\paragraph{Ablations.}
\label{sec:exp-ablations}
\TODO{(i) global variables $\Gv{t}$ on/off --- a measured variable, not an assumption, and the
delta is itself a result since \citet{wu2023probabilistic} propose them but never test them;
(ii) exact vs.\ mean-field readout, below.}

\paragraph{Exact vs.\ mean-field readout.}
\label{sec:exp3}
\TODO{Two variants. (a) Evaluation-time swap: train with the mean-field readout, substitute
\Cref{eq:exact-readout} at evaluation on the same parameters --- isolates the approximation
error, nearly free. (b) Train with each --- measures the end-to-end effect, one extra full run.
Neither outcome is a win or a loss: a small gap says mean field is adequate and cheap, a large
one says the tree structure should be exploited. Also check whether the mixture-of-softmaxes
form relieves the rank-$\nlab$ bottleneck of \Cref{sec:vs-lm-head}.}
